\newpage
\phantomsection
\label{prf:fraction-order-comparison-positive-denominator-product}

\begin{remark*}[Return]
\hyperref[thm:fraction-order-comparison-positive-denominator-product]{Return to Theorem}
\end{remark*}

\begin{theorem*}[Fraction Comparison When the Denominator Product Is Positive]
For any $a,b,c,d\in F$, if $b\neq 0$, $d\neq 0$, and $0<b\cdot d$, then
\[
\frac{a}{b}<\frac{c}{d}
\quad\Longleftrightarrow\quad
a\cdot d<b\cdot c.
\]
\end{theorem*}

\begin{proof}
\textbf{Professional Standard Proof.}~\\
Let $a,b,c,d\in F$ with
\[
b\neq 0,\qquad d\neq 0,\qquad 0<b\cdot d.
\]
By
\hyperref[thm:fraction-order-comparison-general]{Theorem~\ref*{thm:fraction-order-comparison-general}},
we have
\[
\frac{a}{b}<\frac{c}{d}
\quad\Longleftrightarrow\quad
(a\cdot d)\cdot(b\cdot d)<(b\cdot c)\cdot(b\cdot d).
\]
Thus it remains to show that, under the hypothesis $0<b\cdot d$, the latter
inequality is equivalent to
\[
a\cdot d<b\cdot c.
\]

Since $0<b\cdot d$, multiplication by the positive element $b\cdot d$
preserves order.  Therefore, if
\[
a\cdot d<b\cdot c,
\]
then by
\hyperref[thm:multiply-positive]{Theorem~\ref*{thm:multiply-positive}},
\[
(a\cdot d)\cdot(b\cdot d)<(b\cdot c)\cdot(b\cdot d).
\]

Conversely, suppose that
\[
(a\cdot d)\cdot(b\cdot d)<(b\cdot c)\cdot(b\cdot d).
\]
Because $b\cdot d>0$, its inverse is positive by
\hyperref[thm:inverse-positive]{Theorem~\ref*{thm:inverse-positive}}:
\[
(b\cdot d)^{-1}>0.
\]
Applying
\hyperref[thm:multiply-positive]{Theorem~\ref*{thm:multiply-positive}}
to the displayed inequality with positive multiplier $(b\cdot d)^{-1}$, we
obtain
\[
\bigl((a\cdot d)\cdot(b\cdot d)\bigr)(b\cdot d)^{-1}
<
\bigl((b\cdot c)\cdot(b\cdot d)\bigr)(b\cdot d)^{-1}.
\]
Associativity gives
\[
(a\cdot d)\bigl((b\cdot d)(b\cdot d)^{-1}\bigr)
<
(b\cdot c)\bigl((b\cdot d)(b\cdot d)^{-1}\bigr),
\]
so
\[
(a\cdot d)\cdot 1<(b\cdot c)\cdot 1.
\]
Hence
\[
a\cdot d<b\cdot c.
\]

Combining the two equivalences,
\[
\frac{a}{b}<\frac{c}{d}
\quad\Longleftrightarrow\quad
a\cdot d<b\cdot c.
\]
\end{proof}

\begin{remark*}[Proof structure]
The proof starts from the general fraction-comparison theorem, which gives an
equivalent inequality after multiplying both sides by the common factor $bd$.
When $bd$ is positive, that extra factor does not disturb the order: one
direction follows immediately by multiplying by a positive element, and the
other direction is recovered by multiplying by the positive inverse
$(bd)^{-1}$.  So the theorem is the positive-denominator-product simplification
of the general comparison law.
\end{remark*}

\begin{remark*}[Dependencies]~\\
\begin{itemize}
  \item \hyperref[thm:fraction-order-comparison-general]{Theorem~\ref*{thm:fraction-order-comparison-general}}
  \item \hyperref[thm:multiply-positive]{Theorem~\ref*{thm:multiply-positive}}
  \item \hyperref[thm:inverse-positive]{Theorem~\ref*{thm:inverse-positive}}
\end{itemize}
\end{remark*}
